% GENERATED FILE. Edit canonical lessons, then run npm run generate:books.
% canonical-source-hash: fnv1a64:2d97343c6c30f1aa
% canonical-lessons: MR-C282-kasa-jaycha, MR-C282-ujvikade-vala, MR-R282-asking-the-way

\chapter{Asking the Way}
\label{ch:mr-asking-the-way}
\hlchaptermodality{282}

\begin{chapteropening}
\textbf{By the end of this chapter:} \emph{I can ask the way (sṭeśanlā kasaṁ jāycaṁ?) and how far it is, and give the way with saraḷ jā and ujvīkaḍe vaḷā.}

Complete MR-R282-asking-the-way: ask the way to a place and give the way back.
\end{chapteropening}

\section[\texorpdfstring{\mr{स्टेशनला} \mr{कसं} \mr{जायचं}? (sṭeśanlā kasaṁ jāycaṁ?)}{स्टेशनला कसं जायचं? (sṭeśanlā kasaṁ jāycaṁ?)}]{\mr{स्टेशनला} \mr{कसं} \mr{जायचं}? (sṭeśanlā kasaṁ jāycaṁ?) — asking the way}

\label{lesson:MR-C282-kasa-jaycha}



\begin{quote}
Say \textquotedblleft{}where is the station?\textquotedblright{} (\emph{sṭeśan kuṭhe āhe?}) Then \textquotedblleft{}I used to go\textquotedblright{}. (\emph{mī jāyco}.)
\end{quote}

\subsection*{You'll want to know: … \mr{ला} \mr{कसं} \mr{जायचं}?}
\begin{quote}
\textbf{\mr{कृपया}, \mr{स्टेशनला} \mr{कसं} \mr{जायचं}?}
\end{quote}

\begin{quote}
\emph{kṛpayā, sṭeśanlā kasaṁ jāycaṁ?} — \textquotedblleft{}Please, how do I get to the station?\textquotedblright{}
\end{quote}

\textbf{\mr{जायचं}} is the \textbf{-\mr{ायच}-} of \textbf{\mr{जायचो}} with \textbf{-\mr{ं}}: \textquotedblleft{}how is it to be gone\textquotedblright{}. Put \textbf{-\mr{ला}} on the place: \textbf{\mr{बसस्थानकाला}} (\emph{basasthānakālā}), \textbf{\mr{विमानतळाला}} (\emph{vimāntaḷālā}). To ask for a place by name, the question you know stays the best: \textbf{\mr{बसस्थानक} \mr{कुठे} \mr{आहे}?}

\subsection*{You'll want to know: how far}
\begin{quote}
\textbf{\mr{इथून} \mr{किती} \mr{लांब} \mr{आहे}?}
\end{quote}

\begin{quote}
\emph{ithūn kitī lāmb āhe?} — \textquotedblleft{}How far is it from here?\textquotedblright{}
\end{quote}

\textbf{\mr{इथून}} (\emph{ithūn}, \textquotedblleft{}from here\textquotedblright{}) is \textbf{\mr{इथे}} with \textbf{-\mr{ऊन}}, \textquotedblleft{}from\textquotedblright{}. The answers: \textbf{\mr{जवळ} \mr{आहे}} (\emph{javaḷ āhe}), \textquotedblleft{}it is near\textquotedblright{}; \textbf{\mr{खूप} \mr{लांब} \mr{आहे}} (\emph{khūp lāmb āhe}), \textquotedblleft{}it is a long way\textquotedblright{}; \textbf{\mr{दहा} \mr{मिनिटं}} (\emph{dahā miniṭaṁ}), \textquotedblleft{}ten minutes\textquotedblright{}.

\subsection*{Your turn}
Say these aloud:

\begin{itemize}
\raggedright
  \item \textquotedblleft{}Please, how do I get to the station?\textquotedblright{} — \emph{kṛpayā, sṭeśanlā kasaṁ jāycaṁ?}
  \item \textquotedblleft{}How far is it from here?\textquotedblright{} — \emph{ithūn kitī lāmb āhe?}
  \item \textquotedblleft{}Is it near?\textquotedblright{} — \emph{javaḷ āhe kā?}
\end{itemize}

\begin{tcolorbox}[breakable,colback=teal!4,colframe=teal!35!black,title={Before you move on}]
\textquotedblleft{}How do I get to the airport?\textquotedblright{} (\textbf{\mr{विमानतळाला} \mr{कसं} \mr{जायचं}?}) \textquotedblleft{}How far is it from here?\textquotedblright{} (\textbf{\mr{इथून} \mr{किती} \mr{लांब} \mr{आहे}?}) \textquotedblleft{}From here\textquotedblright{}? (\textbf{\mr{इथून}}.)
\end{tcolorbox}
\section[\texorpdfstring{\mr{उजवीकडे} \mr{वळा} (ujvīkaḍe vaḷā)}{उजवीकडे वळा (ujvīkaḍe vaḷā)}]{\mr{उजवीकडे} \mr{वळा} (ujvīkaḍe vaḷā) — giving the way}

\label{lesson:MR-C282-ujvikade-vala}



\begin{quote}
Say \textquotedblleft{}go right\textquotedblright{} with the \textbf{-\mr{कडे}} you know. (\emph{ujvīkaḍe jā}.) Then \textquotedblleft{}straight\textquotedblright{}. (\emph{saraḷ}.)
\end{quote}

\subsection*{You'll want to know: go, turn}
\begin{quote}
\textbf{\mr{सरळ} \mr{जा}, \mr{मग} \mr{चौकात} \mr{उजवीकडे} \mr{वळा}.}
\end{quote}

\begin{quote}
\emph{saraḷ jā, mag caukāt ujvīkaḍe vaḷā.} — \textquotedblleft{}Go straight, then turn right at the square.\textquotedblright{}
\end{quote}

\textbf{\mr{वळा}} is \textbf{\mr{वळणे}} with \textbf{-\mr{आ}}, the polite \textquotedblleft{}turn\textquotedblright{}, built like \textbf{\mr{घ्या}} and \textbf{\mr{बसा}}. To a friend: \textbf{\mr{वळ}} (\emph{vaḷ}). \textbf{\mr{जा}} serves both. With \textbf{\mr{पहिल्या}}, \textquotedblleft{}the first\textquotedblright{}, the way gets precise: \textbf{\mr{पहिल्या} \mr{चौकात} \mr{डावीकडे} \mr{वळा}} (\emph{pahilyā caukāt ḍāvīkaḍe vaḷā}), \textquotedblleft{}turn left at the first square\textquotedblright{}; \textbf{\mr{पुढे} \mr{जा}} (\emph{puḍhe jā}), \textquotedblleft{}go on ahead\textquotedblright{}.

\subsection*{You'll want to know: what stands beside it}
\begin{quote}
\textbf{\mr{स्टेशन} \mr{मंदिराच्या} \mr{समोर} \mr{आहे}.}
\end{quote}

\begin{quote}
\emph{sṭeśan mandirācyā samor āhe.} — \textquotedblleft{}The station is opposite the temple.\textquotedblright{}
\end{quote}

\textbf{-\mr{च्या}} on the landmark, then the place word: \textbf{\mr{दुकानाच्या} \mr{जवळ}} (\emph{dukānācyā javaḷ}), \textquotedblleft{}near the shop\textquotedblright{}; \textbf{\mr{मंदिराच्या} \mr{मागे}} (\emph{mandirācyā māge}), \textquotedblleft{}behind the temple\textquotedblright{}.

\subsection*{Your turn}
Say these aloud:

\begin{itemize}
\raggedright
  \item \textquotedblleft{}Go straight, then turn right at the square\textquotedblright{} — \emph{saraḷ jā, mag caukāt ujvīkaḍe vaḷā}
  \item \textquotedblleft{}Turn left at the first square\textquotedblright{} — \emph{pahilyā caukāt ḍāvīkaḍe vaḷā}
  \item \textquotedblleft{}The station is opposite the temple\textquotedblright{} — \emph{sṭeśan mandirācyā samor āhe}
\end{itemize}

\begin{tcolorbox}[breakable,colback=teal!4,colframe=teal!35!black,title={Before you move on}]
\textquotedblleft{}Turn right\textquotedblright{}? (\textbf{\mr{उजवीकडे} \mr{वळा}}.) \textquotedblleft{}Go straight\textquotedblright{}? (\textbf{\mr{सरळ} \mr{जा}}.) \textquotedblleft{}Opposite the temple\textquotedblright{}? (\textbf{\mr{मंदिराच्या} \mr{समोर}}.)
\end{tcolorbox}
\section[(review)]{Review — asking the way}

\label{lesson:MR-R282-asking-the-way}



\begin{quote}
Say \textquotedblleft{}I must go to the station\textquotedblright{} and \textquotedblleft{}you shouldn't go at night\textquotedblright{}. (\emph{malā sṭeśanlā jāylā pāhije}, \emph{tū rātrī jāū nayes}.) Then say \textbf{\mr{ज्ञान}} and \textbf{\mr{ऋतू}}, and name the letter each one holds: \textbf{\mr{ञ}} and \textbf{\mr{ऋ}}.
\end{quote}

\subsection*{Your turn}
Say these aloud:

\begin{itemize}
\raggedright
  \item Please, how do I get to the airport?
  \item How far is it from here? Is it near?
  \item Go straight, then turn left at the square.
  \item The bus station is near the shop, opposite the temple.
  \item the way from your door to the nearest shop, in three sentences
\end{itemize}

\begin{tcolorbox}[breakable,colback=teal!4,colframe=teal!35!black,title={Before you move on}]
\textquotedblleft{}How do I get to …?\textquotedblright{} (\textbf{… \mr{ला} \mr{कसं} \mr{जायचं}?}) \textquotedblleft{}Turn\textquotedblright{}? (\textbf{\mr{वळा}}.) \textquotedblleft{}Near the shop\textquotedblright{}? (\textbf{\mr{दुकानाच्या} \mr{जवळ}}.) Which letters did \textbf{\mr{ऋतू}} and \textbf{\mr{ज्ञान}} bring? (\textbf{\mr{ऋ}} and \textbf{\mr{ञ}}.)
\end{tcolorbox}
